% gerado por ferramentas/md2tex_enunciado.py a partir de questoes/q094.md (não editar à mão)
Em um reator de fusão nuclear, é necessário produzir plasma a um milhão de graus centígrados para desencadear reações de fusão nuclear controladas. O plasma é formado por um gás ionizado que contém inicialmente deutério $D^{+}$ e trítio $T^{+}$ (isótopos do hidrogênio, com um e dois nêutrons, respectivamente). Nesse plasma, a energia térmica é suficientemente alta para desencadear uma reação de fusão nuclear na qual o núcleo do deutério, partícula $D^{+}$, e o núcleo de trítio, partícula $T^{+}$, colidem para formar uma partícula $\alpha$ e um nêutron \textit{n}. A representação dessa reação de fusão é\par
\[ D^{+} + T^{+} \longrightarrow \alpha + n \]
A partícula $\alpha$ é formada por dois prótons e dois nêutrons. A energia liberada pela fusão é dividida na forma de energia cinética entre a partícula $\alpha$ e o nêutron \textit{n}. Considere as massas do próton e do nêutron iguais, que o momento linear total das partículas reagentes é nulo e também que a energia cinética dos produtos é de 20\,MeV.\par
\fonte{Disponível em: www.scientificamerican.com. Acesso em: 11 nov. 2019 (adaptado).}
Nessa reação de fusão, a energia cinética do nêutron, em megaelétron-volt, é\par
\begin{altlinha}
\alth{A}{4,0\,MeV.}%
\alth{B}{10\,MeV.}%
\alth{C}{13\,MeV.}%
\alth{D}{15\,MeV.}%
\alth{E}{16\,MeV.}
\end{altlinha}
